\documentclass[10pt,a4paper]{amsart}
\usepackage[margin=19mm]{geometry}
\usepackage{amsmath,amssymb,mathtools}
\usepackage[scr=rsfs,cal=euler]{mathalfa}
\usepackage{xfrac}
\usepackage[unicode,hidelinks,bookmarks=false]{hyperref}
\newcommand{\ie}{i.e.\ }
\newcommand{\cf}{cf.\ }
\newcommand{\resp}{respectively\ }
\newcommand{\Subsection}{\subsection}
\providecommand{\nobreakdash}{\nobreak\hbox{-}}
\setlength{\emergencystretch}{2em}
\multlinegap0pt
\usepackage{amssymb}
\usepackage{enumerate}
\newtheorem{theorem}{Theorem}
\title{Shadowing and metric expansivity on Fr\'echet spaces}
\author{Xinxing Wu, Guting Wang}
\date{Source: arXiv:2607.28705v1 (CC BY 4.0)}
\begin{document}
\maketitle

\noindent\emph{Source and licence.} Shadowing and metric expansivity on Fr\'echet spaces. Version arXiv:2607.28705v1 (CC BY 4.0), distributed by arXiv under the Creative Commons Attribution 4.0 International licence, http://creativecommons.org/licenses/by/4.0/, as recorded in the arXiv licence metadata for this exact version and shown on the abstract page https://arxiv.org/abs/2607.28705v1. Full licence notice: \texttt{licenses/arxiv-2607.28705-CC-BY-4.0.txt}.\par\smallskip

\noindent\textit{Editorial scope.} Counted unit: the article's theorem thm:omega-weighted-shift-shadowing (source0.tex lines 683--687). The article's complete original proof of it is delivered with it (source0.tex lines 689--860, one \texttt{proof} environment of 172 lines). No mathematics is written, repaired or re-derived: the statement and the proof are the article's own lines, in the article's own notation, and no retained line is substituted or re-flowed.  Retained uncounted as the source-local context the proof needs: lines 649--652; lines 657--659; lines 672--676.  Nothing was omitted from any retained span, so the package carries no omission record.  No retained line contains a citation, so the wrapper carries no bibliography and no bibliography entry.  No retained line contains a figure environment, an image-inclusion command or drawing code, so no image is delivered and the package declares no asset dependency.  Editorial formatting only, in addition: an AMS \texttt{amsart} preamble that restates the article's own declarations and macros used by the retained lines, namely the environments \texttt{theorem} on separate counters, together with the standard packages those lines need.  The retained environments print in this document as Theorem 1 in order of appearance, whereas the article prints the same items with the numbering of its own full document.  The complete native source of the article as distributed in the arXiv e-print for this exact version is bundled unchanged as \texttt{sources/206434/source0.tex}.
\begin{equation}\label{eq:local_base_0}
\{\{x\in\mathbb{K}^{\mathbb{Z}}:p_m(x)<\varepsilon\}:
m\in\mathbb{N}_0,\ \varepsilon>0\}.
\end{equation}

\begin{equation}\label{eq:Def-bilateral_weighted_forward_shift}
(F_w x)_j=w_{j-1}\cdot x_{j-1}\quad(j\in\mathbb{Z}).
\end{equation}

\begin{equation}\label{eq:Iteration_Formula}
(F_w^n x)_j=W_{j-n}(n)\cdot x_{j-n},
\quad
(F_w^{-n}x)_j=W_j(n)^{-1}\cdot x_{j+n}.
\end{equation}

\begin{theorem}\label{thm:omega-weighted-shift-shadowing}
The bilateral weighted forward shift \(F_w\) defined by
\eqref{eq:Def-bilateral_weighted_forward_shift} with nonzero
weights has the shadowing property.
\end{theorem}

\begin{proof}
For each neighborhood $V_1$ of $0$, by \eqref{eq:local_base_0},
there exist $m\in\mathbb{N}_0$ and $\varepsilon>0$ such that
\[
V:=\{x\in\mathbb{K}^{\mathbb{Z}}:p_m(x)<\varepsilon\}\subset V_1.
\]
Let
\[
C_m^+
=
\max_{1\le j\le m}\sum_{q=1}^{j}|W_q(j-q)|,
\quad
C_m^-
=
\max_{-m\le j\le -1}
|W_j(-j)|^{-1}\cdot \sum_{q=j+1}^{0}|W_q(-q)|,
\]
and
\[
C_m=1+C_m^{+}+C_m^{-},
\quad
\xi=\frac{\varepsilon}{C_m},
\]
where the maximum over an empty set is taken to be \(0\). Choose
\[
U=\{x\in\mathbb{K}^{\mathbb{Z}}:p_m(x)<\xi\}.
\]
For each \(U\)-pseudotrajectory \((x^{(n)})_{n\in\mathbb{Z}}\)
of \(F_w\), put
\begin{equation}\label{eq:Def_e^{n}}
e^{(n)}=x^{(n+1)}-F_{w}x^{(n)}
\quad(n\in\mathbb{Z}).
\end{equation}
This, together with \eqref{eq:Def-bilateral_weighted_forward_shift}, implies
\begin{equation}\label{eq:Def_e_q^{n}}
e_q^{(n)}=x^{(n+1)}_q-(F_{w}x^{(n)})_{q}=
x^{(n+1)}_q-w_{q-1}\cdot x^{(n)}_{q-1}
\quad(n\in\mathbb{Z}).
\end{equation}
By the choice of \(U\), we have
\begin{equation}\label{eq:e_q-Upper-Bound}
|e_q^{(n)}|\leq p_{m}(x^{(n+1)}-F_{w}x^{(n)})<\xi
\quad(n\in\mathbb{Z},\ |q|\le m).
\end{equation}

Choose a sequence \(y\in\mathbb{K}^{\mathbb{Z}}\) as
\[
y_i=
\begin{cases}
W_i(-i)^{-1}\cdot x^{(-i)}_0, & i\le0,\\
W_0(i)\cdot x^{(-i)}_0, & i>0.
\end{cases}
\]
By direct calculation, we have that
\begin{itemize}
  \item For each \(n\geq 0\),
\[
(F_w^n y)_0=W_{-n}(n)\cdot y_{-n}=x^{(n)}_0;
\]
  \item For each \(n<0\),
\[
(F_w^n y)_0=W_0(-n)^{-1}\cdot y_{-n}=x^{(n)}_0;
\]
\end{itemize}
and thus
\begin{equation}\label{eq:F_w^n(y)_0}
(F_w^n y)_0=x^{(n)}_0
\quad(\forall n\in\mathbb{Z}).
\end{equation}

Fix \(n\in\mathbb{Z}\). For each $j\in \mathbb{Z}$ with \(|j|\le m\), to prove
$|(F_w^n y)_j-x^{(n)}_j|<\varepsilon$, we consider the following
three cases:

1)
If \(j>0\), then \(F_w^n y = F_w^j(F_w^{n-j}y)\).
Since \((F_w^{n-j}y)_0 = x^{(n-j)}_0\) by \eqref{eq:F_w^n(y)_0},
and since \((F_w^j x)_j = W_0(j)\cdot x_0\) for any
\(x\in \mathbb{K}^{\mathbb{Z}}\) by \eqref{eq:Iteration_Formula},
we obtain
\begin{equation}\label{eq:F_w^n(y)}
(F_w^n y)_j=W_0(j)\cdot x^{(n-j)}_0.
\end{equation}
Meanwhile, by \eqref{eq:Def_e^{n}}, we have
\begin{align*}
x^{(n)}=& F_{w}x^{(n-1)}+e^{(n-1)}\\
=&
\cdots \\
=& F_{w}^{j}x^{(n-j)}+
F_{w}^{j-1}e^{(n-j)}+\cdots +e^{(n-1)}\\
=& F_{w}^{j}x^{(n-j)}
+\sum_{q=1}^{j}F_{w}^{j-q}e^{(n-j+q-1)},
\end{align*}
and thus, by \eqref{eq:Iteration_Formula}
\[
x^{(n)}_j
=
W_0(j)\cdot x^{(n-j)}_0
+
\sum_{q=1}^{j} W_q(j-q)\cdot e_q^{(n-j+q-1)}.
\]
This, together with \eqref{eq:e_q-Upper-Bound} and \eqref{eq:F_w^n(y)},
implies
\[
\left|(F_w^n y)_j-x^{(n)}_j\right|
=\left|\sum_{q=1}^{j} W_q(j-q)\cdot e_q^{(n-j+q-1)}\right|
\leq \sum_{q=1}^{j}|W_q(j-q)|\cdot \xi
\leq C_{m}^{+}\cdot \xi
< C_{m}\cdot \xi=\varepsilon .
\]

2)
If \(j<0\), by \eqref{eq:Iteration_Formula},
\[
(F_w^{j}x)_j = W_j(-j)^{-1}\cdot x_0
\quad (\forall x\in\mathbb{K}^{\mathbb{Z}}).
\]
By \eqref{eq:F_w^n(y)_0}, taking \(x = F_w^{n-j} y\) gives
\begin{equation}\label{eq:F_w^n(y)-2}
(F_w^{n}y)_j = (F_w^{j}(F_w^{n-j} y))_j = W_j(-j)^{-1}\cdot(F_w^{n-j} y)_0
=W_j(-j)^{-1}\cdot x^{(n-j)}_0.
\end{equation}
Meanwhile, by \eqref{eq:Def_e^{n}}, we have
\begin{align*}
x^{(n-j)}=& F_{w}x^{(n-j-1)}+e^{(n-j-1)}\\
=&
\cdots \\
=& F_{w}^{-j}x^{(n)}+
F_{w}^{-j-1}e^{(n)}+\cdots +e^{(n-j-1)}\\
=& F_{w}^{-j}x^{(n)}
+\sum_{q=j+1}^{0}F_{w}^{-q}e^{(n+q-j-1)},
\end{align*}
and thus, by \eqref{eq:Iteration_Formula}
\[
x^{(n-j)}_0
=
W_j(-j)\cdot x^{(n)}_j
+
\sum_{q=j+1}^{0}W_q(-q)\cdot e_q^{(n+q-j-1)}.
\]
Then
\[
x^{(n)}_j=W_j(-j)^{-1}\cdot x^{(n-j)}_0
-
W_j(-j)^{-1}\cdot \sum_{q=j+1}^{0}W_q(-q)\cdot
e_q^{(n+q-j-1)}.
\]
This, together with \eqref{eq:e_q-Upper-Bound} and \eqref{eq:F_w^n(y)-2},
implies
\begin{align*}
|(F_w^n y)_j-x^{(n)}_j|
=& \left|W_j(-j)^{-1}\cdot \sum_{q=j+1}^{0}W_q(-q)\cdot
e_q^{(n+q-j-1)}\right|\\
\le &
|W_j(-j)|^{-1}\cdot \sum_{q=j+1}^{0}|W_q(-q)|\cdot \xi
\leq C_{m}^{-}\cdot \xi <\varepsilon .
\end{align*}

3) If \(j=0\), by \eqref{eq:F_w^n(y)_0}, then
\[
|(F_w^n y)_0-x^{(n)}_0|=0<\varepsilon.
\]

Thus, $|(F_w^n y)_j-x^{(n)}_j|<\varepsilon$ holds for all \(|j|\le m\).
By the definition of \(p_m\), we have
\[
p_m(F_w^n y-x^{(n)})<\varepsilon
\quad(\forall n\in\mathbb{Z}).
\]
Therefore, \(F_w^n y-x^{(n)}\in V\subset V_1\) for all \(n\in\mathbb{Z}\).
Hence \(F_w\) has the shadowing property.
\end{proof}

\end{document}
